% Problem record op_a381e2ccd80cec9c. This TeX file is derived from
% database/problems_json/op_a381e2ccd80cec9c.json by scripts/sync-tex.mjs; edit the JSON record
% and rerun the script rather than editing this file.
\section{Additivity of the relative entropy of entanglement}
\label{sec:op_a381e2ccd80cec9c}

\paragraph{Problem.}
Does the relative entropy of entanglement of every bipartite state equal its
regularization, or is regularization genuinely necessary?  For a
finite-dimensional bipartite system $A{:}B$, write $\operatorname{Sep}(A{:}B)$
for the set of separable states and
$D(\rho\Vert\sigma)=\operatorname{Tr}[\rho(\log_2\rho-\log_2\sigma)]$ for the
Umegaki relative entropy, defined when
$\operatorname{supp}\rho\subseteq\operatorname{supp}\sigma$, with the trace
evaluated on $\operatorname{supp}\rho$ and $0\log_2 0:=0$.  Define the
relative entropy of entanglement and its regularization by
\begin{equation}
  E_R(\rho)
  :=\min_{\sigma\in\operatorname{Sep}(A{:}B)}D(\rho\Vert\sigma),
  \qquad
  E_R^\infty(\rho)
  :=\lim_{n\to\infty}\frac1nE_R\bigl(\rho^{\otimes n}\bigr).
  \label{eq:rer-definition}
\end{equation}
The limit in Eq.~\eqref{eq:rer-definition} exists and equals
$\inf_{n\geq1}\frac1nE_R(\rho^{\otimes n})$ by Fekete's lemma: the product
of minimizing separable states is separable and $D$ is additive on tensor
products, which gives the subadditivity
$E_R(\rho\otimes\sigma)\leq E_R(\rho)+E_R(\sigma)$, and subadditivity implies
convergence of the normalized terms to their infimum, not that each of them
is nonincreasing.  The archived question is
whether single copies already suffice, that is, whether
\begin{equation}
  E_R^\infty(\rho)=E_R(\rho)
  \quad\text{for every finite-dimensional bipartite state }\rho.
  \label{eq:rer-question}
\end{equation}
Since $E_R^\infty(\rho)\leq E_R(\rho)$ always holds,
Eq.~\eqref{eq:rer-question} can only fail strictly, through a single state
$\rho$ with
\begin{equation}
  E_R^\infty(\rho)<E_R(\rho).
  \label{eq:rer-failure}
\end{equation}

\subsection*{Status}

Solved

\subsection*{Source}

Whether the regularization in Eq.~\eqref{eq:rer-definition} is necessary is
the additivity question raised with the asymptotic relative entropy of
entanglement: Vollbrecht and Werner's study of entanglement measures under
symmetry, which supplied the antisymmetric counterexample recorded in
Progress \sourcecite{ref:rer-vollbrecht-werner}{VW01}, and the closed-form
evaluation by Audenaert et al. of the asymptotic relative entropy of
entanglement taken with respect to the PPT states for Werner states
\sourcecite{ref:rer-audenaert}{AEM+01} together frame the universal equality
of Eq.~\eqref{eq:rer-question}.

\subsection*{Progress}

\begin{itemize}
  \item Vollbrecht and Werner's analysis of entanglement measures under symmetry is
the original source of the non-additivity of $E_R$: the normalized
antisymmetric Werner state
\begin{equation}
  \rho_-:=\frac{2}{d(d-1)}P^{\mathrm{AS}},
  \label{eq:rer-antisymmetric-state}
\end{equation}
on $\mathbb C^d\otimes\mathbb C^d$, with $P^{\mathrm{AS}}$ the projector onto
the antisymmetric subspace, already fails additivity of
Eq.~\eqref{eq:rer-definition} on tensor powers
\sourcecite{ref:rer-vollbrecht-werner}{VW01}.

  \item Audenaert, Eisert, Jan\'e, Plenio, Virmani, and De Moor computed in closed
analytical form the asymptotic relative entropy of entanglement with respect
to the PPT states --- the regularization of
$\min_{\sigma\in\operatorname{PPT}(A{:}B)}D(\rho\Vert\sigma)$ over the set of
states with positive partial transpose --- for Werner states of arbitrary
dimension \sourcecite{ref:rer-audenaert}{AEM+01}, with the companion extension
to orthogonally invariant states
\sourcecite{ref:rer-audenaert-orthogonal}{ADM+02}.  This is a different
quantity from the separable-reference regularization of
Eq.~\eqref{eq:rer-definition}: since
$\operatorname{Sep}(A{:}B)\subseteq\operatorname{PPT}(A{:}B)$, the PPT values
only lower-bound the separable ones, and for the state of
Eq.~\eqref{eq:rer-antisymmetric-state} the PPT-reference value is
$\log_2\frac{d+2}{d}$, which does not determine the separable $E_R^\infty$.

  \item Zhu, Chen, and Hayashi established the exact one- and two-copy values
\begin{equation}
  E_R(\rho_-)=1,
  \qquad
  E_R(\rho_-\otimes\rho_-)=1-\log_2\frac{d-1}{d}<2
  \quad\text{for }d\geq3,
  \label{eq:rer-two-copy-values}
\end{equation}
for the state of Eq.~\eqref{eq:rer-antisymmetric-state}, as instances of
their antisymmetric counterexamples for relative-entropy entanglement
measures \sourcecite{ref:rer-zhu-chen-hayashi}{ZCH10}.  Combined with
subadditivity, Eq.~\eqref{eq:rer-two-copy-values} gives
$E_R^\infty(\rho_-)\leq\frac12\bigl(1-\log_2\frac{d-1}{d}\bigr)<1=E_R(\rho_-)$
for every $d\geq3$, exhibiting Eq.~\eqref{eq:rer-failure} and resolving
Eq.~\eqref{eq:rer-question} negatively; the result is published in
\emph{New Journal of Physics}.

  \item Rubboli and Tomamichel re-derived the values of
Eq.~\eqref{eq:rer-two-copy-values} with the separable optimizer verified
explicitly, extended the antisymmetric counterexample to every entanglement
monotone based on a quantum relative entropy, including the Umegaki $E_R$ of
Eq.~\eqref{eq:rer-definition}, and proved additivity on tensor products
whenever one factor is pure, maximally correlated, GHZ, Bell-diagonal,
isotropic, or generalized Dicke
\sourcecite{ref:rer-rubboli-tomamichel}{RT24}; the result is published in
\emph{Communications in Mathematical Physics}.

  \item The post-resolution frontier treats non-additivity as settled: Beigi,
Rubboli, and Tomamichel characterize, for a broad class of composite
hypothesis-testing and resource-theoretic problems, when regularization of
the Umegaki relative entropy is unnecessary through an if-and-only-if
single-copy optimizer criterion \sourcecite{ref:rer-beigi}{BRT26}, and Lami,
Berta, and Regula prove that the Stein exponent of entanglement testing
equals $E_R^\infty$ while the Sanov exponent equals a single-letter reverse
relative entropy of entanglement \sourcecite{ref:rer-lami}{LBR26}.
\end{itemize}

\subsection*{References}

\begin{enumerate}[label={},leftmargin=4.8em,labelsep=0.5em]
  \footnotesize
  \item[\textup{[VW01]}]\label{ref:rer-vollbrecht-werner}
  K. G. H. Vollbrecht and R. F. Werner,
  ``Entanglement Measures under Symmetry,'' \emph{Physical Review A}
  \textbf{64}, 062307 (2001).
  \href{https://arxiv.org/abs/quant-ph/0010095}{arXiv:quant-ph/0010095}.

  \item[\textup{[AEM+01]}]\label{ref:rer-audenaert}
  K. Audenaert, J. Eisert, E. Jan\'e, M. B. Plenio, S. Virmani, and B. De Moor,
  ``The Asymptotic Relative Entropy of Entanglement,'' \emph{Physical Review
  Letters} \textbf{87}, 217902 (2001).
  \href{https://arxiv.org/abs/quant-ph/0103096}{arXiv:quant-ph/0103096}.

  \item[\textup{[ADM+02]}]\label{ref:rer-audenaert-orthogonal}
  K. Audenaert, B. De Moor, K. G. H. Vollbrecht, and R. F. Werner,
  ``Asymptotic Relative Entropy of Entanglement for Orthogonally Invariant
  States,'' \emph{Physical Review A} \textbf{66}, 032310 (2002).
  \href{https://arxiv.org/abs/quant-ph/0204143}{arXiv:quant-ph/0204143}.

  \item[\textup{[ZCH10]}]\label{ref:rer-zhu-chen-hayashi}
  H. Zhu, L. Chen, and M. Hayashi,
  ``Additivity and Non-Additivity of Multipartite Entanglement Measures,''
  \emph{New Journal of Physics} \textbf{12}, 083002 (2010).
  \href{https://arxiv.org/abs/1002.2511}{arXiv:1002.2511}.

  \item[\textup{[RT24]}]\label{ref:rer-rubboli-tomamichel}
  R. Rubboli and M. Tomamichel,
  ``New Additivity Properties of the Relative Entropy of Entanglement and Its
  Generalizations,'' \emph{Communications in Mathematical Physics}
  \textbf{405}, 162 (2024).
  \href{https://arxiv.org/abs/2211.12804}{arXiv:2211.12804}.

  \item[\textup{[BRT26]}]\label{ref:rer-beigi}
  S. Beigi, R. Rubboli, and M. Tomamichel,
  ``Additivity of Quantum Relative Entropies as a Single-Copy Criterion,''
  \emph{Communications in Mathematical Physics} \textbf{407}, 134 (2026).
  \href{https://arxiv.org/abs/2507.05696}{arXiv:2507.05696}.

  \item[\textup{[LBR26]}]\label{ref:rer-lami}
  L. Lami, M. Berta, and B. Regula,
  ``Asymptotic Quantification of Entanglement with a Single Copy,''
  \emph{Nature Physics} \textbf{22}, 439--445 (2026).
  \href{https://doi.org/10.1038/s41567-026-03182-x}{doi:10.1038/s41567-026-03182-x}.

  \item[\textup{[MI08]}]\label{ref:rer-miranowicz}
  A. Miranowicz and S. Ishizaka,
  ``Closed Formula for the Relative Entropy of Entanglement,''
  \emph{Physical Review A} \textbf{78}, 032310 (2008).
  \href{https://arxiv.org/abs/0805.3134}{arXiv:0805.3134}.
\end{enumerate}

\subsection*{Comment}

The answer to Eq.~\eqref{eq:rer-question} is negative: for every $d\geq3$ the
antisymmetric Werner state of Eq.~\eqref{eq:rer-antisymmetric-state}
satisfies Eq.~\eqref{eq:rer-failure}, by Eq.~\eqref{eq:rer-two-copy-values}
and subadditivity.  The resolving results are peer-reviewed: Vollbrecht and
Werner in \emph{Physical Review A}
\sourcecite{ref:rer-vollbrecht-werner}{VW01}, Zhu, Chen, and Hayashi in
\emph{New Journal of Physics} \sourcecite{ref:rer-zhu-chen-hayashi}{ZCH10},
and Rubboli and Tomamichel in \emph{Communications in Mathematical Physics},
which verified the separable optimizer and extended the counterexample to
every relative-entropy-based monotone
\sourcecite{ref:rer-rubboli-tomamichel}{RT24}.  The reverse inequality
$E_R^\infty(\rho)>E_R(\rho)$ cannot occur, since subadditivity is trivial.  The closed forms of
\sourcecite{ref:rer-audenaert}{AEM+01} and
\sourcecite{ref:rer-audenaert-orthogonal}{ADM+02} concern the PPT-reference
regularization, a different quantity, and do not by themselves evaluate the
separable regularization of Eq.~\eqref{eq:rer-definition}.
Two residues lie outside this record: deciding additivity for a given state
is governed by the single-copy optimizer criterion of
\sourcecite{ref:rer-beigi}{BRT26}, and the general two-qubit case turns on
the closed formula and explicit optimizer archived in the two-qubit
relative-entropy-of-entanglement record, weak additivity being known there
for states commuting with their optimizer
\sourcecite{ref:rer-miranowicz}{MI08}; the entanglement-cost record for a
qubit Bell-diagonal state already uses the Bell-diagonal evaluation
$E_R^\infty=L(p_*)$ \sourcecite{ref:rer-zhu-chen-hayashi}{ZCH10}.

\subsection*{Field}

Quantum Resource Theory

\subsection*{Topic}

Entanglement measures; Additivity and regularization; Quantum relative entropy

\subsection*{ID}

\texttt{op\_a381e2ccd80cec9c}
